\chapter{Conclusão}
\label{cap:conclusao}

% ─── TODO: escrever após conclusão dos resultados ─────────────────────────────

Este trabalho apresentou o projeto e a implementação da
\textit{Rizoma}, uma plataforma web acadêmica para análise
integrada de micobioma e metataxonômica, desenvolvida para suportar
três projetos científicos ativos na UFVJM: INOVAHERB, Pós-Fogo e
Biorremediação.

% Principais contribuições:
% - Arquitetura sem Redis (PG LISTEN/NOTIFY)
% - R Worker desacoplado com output 100% estruturado
% - Autenticação institucional Google OAuth
% - Painel unificado de 6 PCoAs como produto científico

\section{Limitações}

% Listar as limitações identificadas (tech debt do CLAUDE.md):
% - WebSocket não conectado ao NOTIFY
% - GSEA hardcoded para org.Hs.eg.db
% - Worker single-threaded
% - CORS restrito a localhost:3000

\section{Trabalhos Futuros}

\begin{itemize}
  \item Integrar o WebSocket \texttt{/jobs/ws/status} ao
        \texttt{LISTEN/NOTIFY} do PostgreSQL para atualizações em
        tempo real sem \textit{polling}.
  \item Parametrizar o banco de anotações do GSEA para suportar
        outros organismos além de \textit{Homo sapiens}.
  \item Implementar paralelismo no R Worker via processos filhos
        para reduzir o bloqueio da fila em análises longas (SpiecEasi,
        Random Forest).
  \item Integrar o PICRUSt2 diretamente ao pipeline, eliminando
        a dependência de arquivo local pré-gerado.
  \item Publicar a plataforma em endereço acessível fora do campus
        via configuração de CORS e exposição do Ingress k3s.
\end{itemize}
